\section[第3章\q 你比自己想像的更富有]{第3章\\ 你比自己想像的更富有}
\label{sec:ch3}

“那就做点什么吧！”芭芭拉如是说。次日早上，柔伊还真的有了行动。她比平常早15分钟做好出门准备。她看不出和咖啡师聊天的必要，但既然芭芭拉都提议了，她还是可以在咖啡厅里多消磨一点时光的，至少可以贴近看看那张照片。

她点了餐，排队等餐，拿到特浓拿铁后，在店里晃了晃。店的内墙是裸砖式的，拱形屋顶漆成了全黑色，大型的全光谱灯泡悬在空中，让咖啡厅内看起来像是布鲁克林区的艺术展厅。很艺术，但有点老派。

她沿着咖啡厅长廊走了一圈，停留在挂照片的角落。有些是漂亮的风景：大雪覆盖的山脉、站在激流中取景的澎湃河流、景深的森林小径。有些地方，因为她在旅行杂志工作，所以能认得出来。有一张是中国长城，另一张是意大利皮尔蒙特区的葡萄庄园，几名年轻人正在葡萄园里劳作，还有一张是秘鲁热带雨林区羽色艳丽的金刚鹦鹉。

张张都漂亮，但她的脚步却没有停下来，直到走到那张照片前方。

就是这张，没错。她退后几步，细细盯着照片看。

说来，那并不是一幅非常壮观的照片，至少表面看来如此。一个晨光中的海滨小村，一条小渔船位于画面右侧，正要出海，还有往来的村民，过着他们日常的渔村生活。

为什么这样勾住她的心呢？

她走近几步，看到照片下方印刷的小字。喔，那是标价：1,200美元。

柔伊的心一沉。价格不斐啊，不过，也算合理，不是吗？再说，这个数字也不是太高，只比她的房租略低一点，柔伊不是负担不起。但她已经很久没看到银行帐面上有一千多元的余额够她随意花费。

她想了一下，确定她从来没有这样一笔可以随意花用的钱。

她弯了腰，再次打量标签。想看一看照片是在哪里拍的，但标签上没有。除了标价，唯二的文本就是照片的标题，标题非常简洁：

“来吧！”

“来吧！”对一幅海滨小村的图片，这是一个看似一点也不搭的标题。来做什么？尽管现在再看这幅画时，柔伊觉得这个标题对她来说也不是那样不搭调。这是哪？应该是希腊的一个小岛屿？“你们在哪儿？”她自言自语道：“罗德岛？圣托里尼岛？”不，应该不是。“克里特岛？”

\subsection{买得起拿铁，就买得起画}

“米克诺斯岛。”

一位先生的声音离得这样近，柔伊几乎惊跳起来，手上的拿铁也差点打翻。

“对不起，”这位先生说：“无意要吓到你，你看得太专注了。”他朝照片示意了一下。“喜欢这张照片？”

柔伊点头。“风景好棒，光线也刚好，叫人想去。”她点头，指了指针题。年长男士顺着她手指的方向看了看，点头同意。她伸出手来：“我是柔伊。柔伊•丹尼斯。”

年长男士握了握她的手，他的手又干又凉，好像一块细致的帆布。“亨利•海顿”，他说：“跟作曲家海顿同姓。”他又补充了一句：“不过，我不是名人。”

“亨利，你好。”现在她当然认出来了，亨利就是咖啡师。“你比自己料想的更有名喔！”

咖啡师低下头来，仿佛在说，是吗？

“我老板对我说起你了。”柔伊解释说：“她说我应该进来，并找你谈一谈。”

“喔，”他回应：“谈什么？”

柔伊开口想说话，但是随即又闭上了，最后她朝亨利撇嘴说：“其实，我不知道要谈什么。”

他笑了笑，对着照片的方向点头示意。“少有人特别钟情这张照片的，”他说：“多数人更喜欢更生动的景象，你知道，就是高山呀，峡谷啊，激流之类的。”

柔伊懂得亨利说的意思。“这张照片，”她说：“我觉得很有生气。”

亨利点头。“就我个人来说，也最喜欢这一张。”

柔伊站直身体，转了360度，把咖啡厅环视一周，对亨利说：“我也是。”

亨利再次低下头。“对了，顺便说一下，这张尚未卖出。”

柔伊笑道：“我希望可以买下来。但是我怕负担不起。”

亨利对着她手中的拿铁示意一下。“如果你买得起拿铁，”他说，又示意了一下墙上的画，“你就买得起这张画。”

“什么？”柔伊问。她有没有听错？这话看似说不通啊！

“或许，”亨利说。“你比自己想像的更富有。”

柔伊不解的笑了。她想，这个说法倒也新奇，但她喜欢亨利的正能量。“这个想法不错。”她说：“但是说真的，我只是看一看。”柔伊又趋身靠近了看照片上的细节：鹅卵石铺出来的小路、白漆房子，还有宝蓝的门和窗台。“你认为是米克诺斯岛？”

亨利也趋近了看，缓缓点头说：“我确定。”

“真是漂亮。”柔伊赞叹说：“我真正想要的是，”她轻声说，好像在自言自语。“真的过去那里，呼吸带着海水咸味的空气，听一听海鸥的声音。用自己的眼睛和耳朵来感受照片中的情景。”

柔伊再次直起身体，不好意思的笑了。用正常的声音说道：“不管怎么说，那都是完全不可能的事。”

“完、全、不、可、能。”亨利重复了一遍，一字一顿的说，好像在玩味文本后面的意思。他又朝柔伊低下头来：“那要看你如何看待能不能这件事了，不是吗？”

柔伊不知道该怎么回答。

“看来你喜欢摄影。”他说：“那你告诉我，你是如何理解‘纽约之眼’的？”

“你说在富尔顿转运站那里吗？”柔伊回答。“我现在正要走过去。”

“不不不，”他回答说：“我不是说那个建筑。我的意思是，从摄影学的角度，如何看待眼窗（oculus）这个词。”

柔伊皱起了眉头。

“眼窗，”他重复一遍，“眼睛如同窗户，摄影镜头也如同窗户。你站在哪里，从站着地方看到什么，是拍出漂亮照片的关键。因为那取决了你想要看到什么。你了解我说的吗？”

柔伊点头，尽管其实并不十分了解亨利的话。

“从摄影角度来说，”咖啡师继续说，“眼窗就是你安置镜头的地方。眼窗的英文取自拉丁文的眼睛之意，镜头的确就是眼睛，因为其实是你先看到景象的，在自己的脑中先看到景象。在你自己的眼中。”

“喔！”柔伊说。从建物底下走过时，她从未想过这个词的意思。

“现在我说的是摄影，”他补充，“说故事其实也是同样的道理。一趟即将开始的旅程也一样。要为即将上门的朋友，准备餐点也是一样。重点是，你站在那里，就有了三样东西：一个是你，一个是世界，另一个是镜头。你想要创造一点什么？”

芭芭拉怎么说他来的？他是一座宝山。但对柔伊来说，跃入脑中的词汇却是“古怪”。但是也不能说他不风度翩翩，不正面积极。明显的老式风格，就像这家咖啡厅的风格。

亨利朝咖啡厅柜台看了看，以确认那边不需要他帮忙。柜台后方，一个戴着圆帽、蓄着胡髭的布鲁克林型男留意到亨利的眼光，他大声说：“别担心，亨利。我们可以的。”

亨利回头看着柔伊，朝角落的高脚桌指了指。“坐下，小聊一下？”

柔伊笑着说：“当然好。”

她随着亨利走到小桌子旁边，分别坐上了两张高脚椅。亨利把他那本欧洲知识分子最爱的意大利Moleskine笔记本翻开来，从封面内侧口袋抽出素描笔，开始在纸片上涂画着什么，只见他手指在笔记本上飞舞。几分钟后，他把笔记本递给柔伊看。

笔记本内页描摹了一块碑石，上头写着：

{\kaishu 柔伊•丹尼尔}

{\kaishu 生于？？ 卒于？？}

“我们假设，这就是你的人生好了。”

“真的，”柔伊淡淡的说：“太可怜了，这么早就死了。”

亨利笑起来，“你还取笑我。我们假设现在要来写一个墓志铭。这个就是你的眼窗，”他用笔敲着素描，“这是你现在站立的地方，回头展望过去：也就是你的人生。那么，这会是怎样的风景呢？”

柔伊凝住呼吸。

她一下无法用语言表达，但是亨利刚刚说的，就是过去几天困扰她的。她现在的人生风景究竟如何？她不知道。

{\kaishu 如果你不知道要去哪里，}

{\kaishu 只怕去了哪都不会快乐。}

“你明白了吗？”亨利说：“图景首先出现在脑中，然后才会实现。所有的风景都是先出现在脑中。那个风景会召唤你，那是你的眼窗。”

柔伊的手机响了，她瞄了一眼，是一位早班实习生传来的，想知道要从哪张图开始编修。

“你要回去工作了？”亨利主动问道。

“我真的要先离开了。”柔伊语带歉意：“非常感谢……有这个机会闲聊一下。”她不知道该称这个简短的谈话为什么，一堂关于如何取景的艺术课？

“这场谈话很愉快，”就在柔伊起身，准备往门口走过去的时候，亨利说：“随时欢迎继续聊下去。”

\subsection{花的比赚的多，人生却缺少了生活}

柔伊来到33楼时，编辑部已经开始全速运转了。她和实习生有个3分钟的短会，交代了代办事项，又和美术组的开了一个小会，然后就打开笔电，被工作淹没了。

尽管如此，她也不时在潜意识想着和海伦娜咖啡师的闲聊。芭芭拉怎么说他来的？他看事情的角度很不一样。“这个一点都没错。”柔伊自言自语道。简短的闲聊，愈想愈觉得很不一样。

{\kaishu 你站在哪里，从站着地方看到什么，是拍出漂亮照片的关键。因为那取决了你想要看到什么。你了解我说的吗？}

坦白说，一点都不了解。

对了，他还提到了她的咖啡。

{\kaishu 如果你买得起拿铁，你就买得起那幅画。}

还有这句：

{\kaishu 或许，你比自己想像更富有。}

都在说些什么呢？

这个晚上，柔伊没有睡好。

事实上，她每晚都没睡好。通常，她总在凌晨两、三点之间醒过来，然后就清醒地躺着，回不到睡梦中。她也没有为什么特别的事情焦虑，但就是说不上原因的焦虑着。

但是这个晚上比其他夜晚更严重一点。因为，凌晨醒来之后，她又昏昏睡去了，而焦虑伴随她进入梦中。

梦中，她在健身房的跑步机上。突然，机器加速了，但她并没有按加速键。没问题，她加快脚步就是。可是，机器又再次加速了。她开始全速奔跑。她开始按减速键，但是却愈按愈快。她现在跑得上气不接下气，心都快要跳出来了，依然赶不上跑步机的速度。

她在气喘吁吁中醒过来，T恤已经被汗浸湿了。她在黑暗中摸到放在床头柜上的水，慢慢地，她的眼睛适应了黑暗，她的心跳也从怦怦跳，回复到正常。

不需要有心理学的学位，谁都能知道这个梦的意思。她清醒时就是在一部每周运转50个小时的跑步机上，这台跑步机还不受她控制。早上，从布鲁克林到曼哈顿，晚上从曼哈顿回到布鲁克林，赚钱，花钱。通常，花的比赚的多。用一个残酷的比喻，就是她用自己的生命全速奔跑，却不知道要去哪里。

在半明半暗的光线中，她打量着公寓四壁，总在这样的时刻，特别诚实的面对自己，她感觉到自己的人生中缺了一点什么，而缺了的那一点，又是非常重要的。那是什么？是爱吗？她还年轻，还有大把的时间。朋友吗？她有洁西卡，还有其他好些朋友。

到底缺了什么呢？她想，应该是缺少生活。
